\section{Week 9: Bayesian Modeling I: Priors, Posteriors and Grid
Approximation}

\begin{quote}
\textbf{Chapter 4: Bayesian modeling} · Lecture 1 h · Lab 2 h
\textbf{Running data:} a recall proportion (30 of 50) and the Week 2 SDT
participant (68/100 hits, 11/100 false alarms).
\end{quote}

\begin{center}\rule{0.5\linewidth}{0.5pt}\end{center}

\subsection{Key terms}

Important terms of this lesson, with their Persian equivalents (repeated
even if they appeared in earlier weeks).

\begin{longtable}[]{@{}
  >{\raggedright\arraybackslash}p{(\linewidth - 4\tabcolsep) * \real{0.3333}}
  >{\raggedright\arraybackslash}p{(\linewidth - 4\tabcolsep) * \real{0.3333}}
  >{\raggedright\arraybackslash}p{(\linewidth - 4\tabcolsep) * \real{0.3333}}@{}}
\caption{\centering Key terms}\tabularnewline
\toprule\noalign{}
\begin{minipage}[b]{\linewidth}\raggedright
English
\end{minipage} & \begin{minipage}[b]{\linewidth}\raggedright
\RL{\CMfa فارسی}
\end{minipage} & \begin{minipage}[b]{\linewidth}\raggedright
\RL{\CMfa توضیح کوتاه}
\end{minipage} \\
\midrule\noalign{}
\endfirsthead
\toprule\noalign{}
\begin{minipage}[b]{\linewidth}\raggedright
English
\end{minipage} & \begin{minipage}[b]{\linewidth}\raggedright
\RL{\CMfa فارسی}
\end{minipage} & \begin{minipage}[b]{\linewidth}\raggedright
\RL{\CMfa توضیح کوتاه}
\end{minipage} \\
\midrule\noalign{}
\endhead
\bottomrule\noalign{}
\endlastfoot
prior & \RL{\CMfa توزیع پیشین} (prior) & \RL{\CMfa باور درباره‌ی پارامتر پیش از دیدن
داده} \\
posterior & \RL{\CMfa توزیع پسین} (posterior) & \RL{\CMfa باور درباره‌ی پارامتر بعد
از دیدن داده} \\
grid approximation & \RL{\CMfa تقریب شبکه‌ای} (grid) & \RL{\CMfa محاسبه‌ی پسین روی
شبکه‌ای از مقادیر} \\
signal detection theory (SDT) & \RL{\CMfa نظریه‌ی تشخیص علامت} & \RL{\CMfa مدلی که
حساسیت را از سوگیری پاسخ جدا می‌کند} \\
hit & \RL{\CMfa اصابت} (hit) & \RL{\CMfa پاسخ «بله» وقتی علامت وجود دارد} \\
Bayes' rule & \RL{\CMfa قاعده‌ی بیز} & \RL{\CMfa پسین} \ensuremath{\propto} \RL{\CMfa درست‌نمایی} ×
\RL{\CMfa پیشین} \\
parameter & \RL{\CMfa پارامتر} & \RL{\CMfa کمیت قابل تنظیم مدل که بین افراد یا شرایط
فرق می‌کند} \\
likelihood & \RL{\CMfa درست‌نمایی} & \RL{\CMfa احتمال داده به‌عنوان تابعی از
پارامترها} \\
evidence & \RL{\CMfa شواهد} & \RL{\CMfa اطلاعاتی که برای تصمیم جمع می‌شود} \\
beta-binomial model & \RL{\CMfa مدل} beta-binomial & \RL{\CMfa پیشین} Beta \RL{\CMfa با
داده‌ی دوجمله‌ای} \\
binomial distribution & \RL{\CMfa توزیع دوجمله‌ای} & \RL{\CMfa توزیع تعداد موفقیت در}
n \RL{\CMfa کوشش} \\
MAP estimate & \RL{\CMfa برآورد} MAP & \RL{\CMfa مُد توزیع پسین} \\
credible interval / HDI & \RL{\CMfa بازه‌ی باورپذیر} (credible interval) &
\RL{\CMfa بازه‌ای که با احتمال مشخص پارامتر را دارد} \\
confidence interval (CI) & \RL{\CMfa بازه‌ی اطمینان} & \RL{\CMfa بازه‌ای که با روش
مشخص پارامتر واقعی را در بر می‌گیرد} \\
prior predictive check & \RL{\CMfa بررسی پیش‌بینی پیشین} & \RL{\CMfa شبیه‌سازی داده از
پیشین برای سنجش معقول بودن آن} \\
simulation & \RL{\CMfa شبیه‌سازی} & \RL{\CMfa تولید داده‌ی مصنوعی با اجرای مدل} \\
sensitivity (d′) & \RL{\CMfa حساسیت} & \RL{\CMfa فاصله‌ی توزیع علامت و نویز} \\
estimate / estimator & \RL{\CMfa برآورد} / \RL{\CMfa برآوردگر} & \RL{\CMfa مقدار به‌دست‌آمده
برای پارامتر} / \RL{\CMfa روش به دست آوردن آن} \\
maximum likelihood estimate (MLE) & \RL{\CMfa برآورد بیشینه‌ی درست‌نمایی} &
\RL{\CMfa پارامترهایی که داده را محتمل‌ترین می‌کنند} \\
cognitive model & \RL{\CMfa مدل شناختی} & \RL{\CMfa توصیف دقیق و ریاضی یا محاسباتی
از یک فرایند ذهنی} \\
trial & \RL{\CMfa کوشش} (trial) & \RL{\CMfa یک بار انجام تکلیف توسط شرکت‌کننده} \\
conjugate prior & \RL{\CMfa پیشین مزدوج} (conjugate) & \RL{\CMfa پیشینی که پسین
هم‌خانواده می‌دهد} \\
forgetting curve & \RL{\CMfa منحنی فراموشی} & \RL{\CMfa کاهش یادآوری با گذشت
زمان} \\
accuracy & \RL{\CMfa دقت} & \RL{\CMfa نسبت پاسخ‌های درست} \\
\end{longtable}

\begin{center}\rule{0.5\linewidth}{0.5pt}\end{center}

\subsection{Learning outcomes}

By the end of this week, students will be able to:

\begin{enumerate}
\def\labelenumi{\arabic{enumi}.}
\tightlist
\item
  \textbf{State} Bayes' rule for parameters and \textbf{identify} the
  roles of prior, likelihood, evidence and posterior.
\item
  \textbf{Compute} posteriors analytically for the beta-binomial model
  and numerically by grid approximation, for one and two parameters.
\item
  \textbf{Summarize} a posterior (mean, median, MAP, SD, equal-tailed
  and highest-density intervals) and \textbf{distinguish} a credible
  interval from a confidence interval.
\item
  \textbf{Explain} how the prior and the amount of data jointly
  determine the posterior, and \textbf{use} prior predictive simulation
  to check whether a prior is sensible.
\item
  \textbf{Interpret} posterior probabilities of psychologically
  meaningful hypotheses (e.g., ``\(P(d' > 1 \mid \text{data})\)'').
\end{enumerate}

\begin{center}\rule{0.5\linewidth}{0.5pt}\end{center}

\subsection{Motivation \& big picture}

Since Week 3 we have asked: ``which single parameter value best accounts
for the data?'' and then, in Week 5, ``how much would that estimate vary
across imaginary replications?'' Researchers often want to say something
more direct:

\begin{quote}
\emph{``Given these data, there is a 95\% probability that this
participant's sensitivity lies between 1.2 and 2.1.''} \emph{``The
probability that the payoff manipulation made participants more liberal
is 0.98.''}
\end{quote}

Frequentist confidence intervals do \textbf{not} license such statements
(Week 5). Bayesian inference does, at a price: we must state what we
believed about the parameters \textbf{before} seeing the data, in the
form of a \textbf{prior distribution}. Bayes' rule then combines the
prior with the likelihood (which we already know how to write) into a
\textbf{posterior distribution}: everything the model and the data
together say about the parameters.

Bayesian methods also make several of our earlier headaches easier:
uncertainty comes for free (no Hessians or bootstraps), hierarchical
models of many participants become natural (Week 13), and model
comparison gets a principled answer (Week 10). The cost is computation,
which is why Weeks 11--13 are about sampling.

\begin{quote}
\textbf{Reminder (Week 1):} this week is about Bayesian \textbf{data
analysis}, a method we as scientists use to estimate the parameters of
\emph{any} model. It makes no claim that the mind is Bayesian. Bayesian
\emph{models of cognition} come in Week 15.
\end{quote}

\paragraph{Where is this used?}

\begin{itemize}
\tightlist
\item
  \textbf{Small samples and individual cases.} Clinical neuropsychology
  often has one patient and a few dozen trials. A Bayesian analysis with
  a sensible prior gives usable, honest uncertainty where MLE is
  unstable.
\item
  \textbf{Direct probability statements} that clinicians and policy
  makers understand: ``there is a 97\% probability that the new training
  improved memory''.
\item
  \textbf{Adaptive testing}: computerized tests (e.g.~vision and hearing
  tests, adaptive exams) update a posterior after every response and
  choose the next item where it is most informative.
\item
  \textbf{Combining evidence}: today's posterior is tomorrow's prior, so
  results of earlier studies can be built in formally.
\end{itemize}

\begin{center}\rule{0.5\linewidth}{0.5pt}\end{center}

\Needspace{11\baselineskip}

\subsection{Lecture notes}

\paragraph{In plain words}

So far we found \emph{one} best parameter value and then asked how much
it might vary. The Bayesian approach instead describes our knowledge
about the parameter as a whole \textbf{probability distribution}. Before
the data, the \textbf{prior} says which values we consider plausible
(e.g.~``recall probability is probably moderate, not 0 or 1''). The data
then re-weight every possible value by its likelihood (Week 4). The
result, the \textbf{posterior}, says how plausible each value is
\emph{after} seeing the data. In words: \emph{posterior \ensuremath{\propto} likelihood ×
prior}. From the posterior we can read off a best guess (its mean), an
uncertainty (its spread), and direct statements such as ``the
probability that \(d' > 1\) is 0.93''. With little data the prior
matters a lot; with much data the posterior is dominated by the
likelihood and ends up close to the MLE. Computing the posterior is easy
for one or two parameters (on a grid) and needs sampling methods beyond
that (Weeks 11--13).

\textbf{Suggested timing (60 min):} Bayes' rule 10 · priors 10 ·
beta-binomial 12 · grid approximation 10 · summaries and credible
intervals 12 · prior predictive checks 6.

\subsubsection{1. Bayes' rule for parameters}

\[
\underbrace{p(\boldsymbol\theta \mid \mathbf{y})}_{\text{posterior}}
= \frac{\overbrace{p(\mathbf{y} \mid \boldsymbol\theta)}^{\text{likelihood}}\;\overbrace{p(\boldsymbol\theta)}^{\text{prior}}}{\underbrace{p(\mathbf{y})}_{\text{evidence (marginal likelihood)}}},
\qquad
p(\mathbf{y}) = \int p(\mathbf{y}\mid\boldsymbol\theta)\,p(\boldsymbol\theta)\,d\boldsymbol\theta.
\]

(We suppress ``\(\mathcal{M}\)'' in the notation this week: everything
is conditional on the model.)

\begin{itemize}
\tightlist
\item
  The \textbf{likelihood} is exactly the function from Week 4.
\item
  The \textbf{prior} expresses what is known or plausible before the
  data.
\item
  The \textbf{evidence} \(p(\mathbf{y})\) does not depend on
  \(\boldsymbol\theta\). For estimation it is just the normalizing
  constant that makes the posterior integrate to 1:
\end{itemize}

\[
p(\boldsymbol\theta \mid \mathbf{y}) \propto p(\mathbf{y}\mid\boldsymbol\theta)\, p(\boldsymbol\theta).
\]

\textbf{Posterior \ensuremath{\propto} likelihood × prior.} In Week 10 the evidence becomes
the star of Bayesian model comparison.

\textbf{Intuition.} The posterior is a compromise: it is high where the
parameter values are both plausible \emph{a priori} and make the data
probable. With little data the prior matters; with a lot of data the
likelihood dominates and the posterior concentrates around the MLE.

\paragraph{Reading the math}

\begin{itemize}
\tightlist
\item
  \(p(\boldsymbol\theta \mid \mathbf{y})\): ``probability (density) of
  the parameter value \(\boldsymbol\theta\) \emph{given} the data''.
  This is what we want.
\item
  \(p(\mathbf{y} \mid \boldsymbol\theta)\): the likelihood, read as a
  function of \(\boldsymbol\theta\) (Week 4).
\item
  \(p(\boldsymbol\theta)\): the prior, chosen before the data.
\item
  \(p(\mathbf{y}) = \int p(\mathbf{y}\mid\boldsymbol\theta)p(\boldsymbol\theta)\,d\boldsymbol\theta\):
  the integral adds up ``likelihood × prior'' over \emph{all} possible
  parameter values. It is one number, the same for every
  \(\boldsymbol\theta\), so for estimation it only rescales the curve to
  total area 1.
\item
  \(\propto\) means ``proportional to'': equal up to a constant factor.
  In practice we compute likelihood × prior at many values and then
  divide by their sum, exactly as in the 4-point table of the worked
  example (0.0004479 / 0.0009574 = 0.468).
\end{itemize}

\subsubsection{2. Choosing priors}

\begin{longtable}[]{@{}
  >{\raggedright\arraybackslash}p{(\linewidth - 4\tabcolsep) * \real{0.3333}}
  >{\raggedright\arraybackslash}p{(\linewidth - 4\tabcolsep) * \real{0.3333}}
  >{\raggedright\arraybackslash}p{(\linewidth - 4\tabcolsep) * \real{0.3333}}@{}}
\caption{\centering 2. Choosing priors}\tabularnewline
\toprule\noalign{}
\begin{minipage}[b]{\linewidth}\raggedright
Kind
\end{minipage} & \begin{minipage}[b]{\linewidth}\raggedright
Example for a recall probability \(\theta\)
\end{minipage} & \begin{minipage}[b]{\linewidth}\raggedright
When
\end{minipage} \\
\midrule\noalign{}
\endfirsthead
\toprule\noalign{}
\begin{minipage}[b]{\linewidth}\raggedright
Kind
\end{minipage} & \begin{minipage}[b]{\linewidth}\raggedright
Example for a recall probability \(\theta\)
\end{minipage} & \begin{minipage}[b]{\linewidth}\raggedright
When
\end{minipage} \\
\midrule\noalign{}
\endhead
\bottomrule\noalign{}
\endlastfoot
Flat / uniform & \(\mathrm{Beta}(1, 1)\) & When every value in the range
is genuinely equally plausible (rarely true) \\
Weakly informative & \(\mathrm{Beta}(2, 2)\): mild preference for
moderate values & \textbf{Default for most cognitive models}: rules out
absurd values, otherwise lets the data speak \\
Informative & \(\mathrm{Beta}(20, 5)\): ``recall after one day is around
0.8'' & When previous studies or theory provide real knowledge \\
\end{longtable}

Principles: - Priors must be fixed \textbf{before} seeing the data. -
They should respect the parameter's range (e.g., no prior mass on
negative decay rates). - They should be checked by \textbf{prior
predictive simulation} (§6): which \emph{data} does the prior imply? - A
``vague'' prior on one scale can be very informative on another scale
(lab Part D). - Report your priors and check the sensitivity of
conclusions to reasonable alternatives (Week 10).

\subsubsection{3. A closed-form case: the beta-binomial model}

For \(y\) successes in \(n\) binomial trials with a
\(\mathrm{Beta}(a_0, b_0)\) prior on \(\theta\):

\[
p(\theta \mid y) \propto \underbrace{\theta^{y}(1-\theta)^{n-y}}_{\text{likelihood}} \cdot \underbrace{\theta^{a_0 - 1}(1-\theta)^{b_0 - 1}}_{\text{prior}}
= \theta^{a_0 + y - 1}(1-\theta)^{b_0 + n - y - 1},
\]

so the posterior is again a Beta distribution:

\[
\theta \mid y \sim \mathrm{Beta}(a_0 + y,\; b_0 + n - y).
\]

When prior and posterior come from the same family, the prior is called
\textbf{conjugate}.

\begin{itemize}
\tightlist
\item
  \(a_0\) and \(b_0\) act like \textbf{pseudo-counts} of prior successes
  and failures.
\item
  Posterior mean: \(\dfrac{a_0 + y}{a_0 + b_0 + n}\), a weighted average
  of the prior mean \(\frac{a_0}{a_0 + b_0}\) and the observed
  proportion \(y/n\), with weights \(a_0 + b_0\) and \(n\).
\item
  \textbf{Connection to Weeks 2 and 4:} the log-linear correction
  \((y + 0.5)/(n + 1)\) is exactly the posterior mean under a
  \(\mathrm{Beta}(0.5, 0.5)\) prior, the so-called \emph{Jeffreys
  prior}. The ``correction'' was a Bayesian estimate in disguise.
\end{itemize}

\paragraph{Reading the math}

\begin{itemize}
\tightlist
\item
  The \textbf{Beta distribution} \(\mathrm{Beta}(a, b)\) is a flexible
  bell-like curve on the interval 0--1, ideal for describing beliefs
  about a probability. Its mean is \(a/(a+b)\). \(\mathrm{Beta}(1,1)\)
  is flat; \(\mathrm{Beta}(20, 5)\) is a narrow bump around 0.8.
\item
  The update rule
  \(\mathrm{Beta}(a_0, b_0) \to \mathrm{Beta}(a_0 + y, b_0 + n - y)\) is
  just \textbf{counting}: add the observed successes to \(a\) and the
  failures to \(b\). The prior behaves as if we had already seen \(a_0\)
  successes and \(b_0\) failures (pseudo-counts).
\item
  \textbf{Example:} prior \(\mathrm{Beta}(2, 2)\) (``as if'' 2 recalls
  and 2 failures), data 7 of 10 \ensuremath{\rightarrow} \(\mathrm{Beta}(9, 5)\), mean
  \(9/14 = 0.643\). The prior pulls the observed 0.70 slightly toward
  0.5, because it adds 4 pseudo-observations to 10 real ones.
\end{itemize}

\subsubsection{4. Grid approximation}

Most cognitive models have no conjugate prior. If there are only one or
two parameters, we can approximate the posterior on a grid:

\begin{enumerate}
\def\labelenumi{\arabic{enumi}.}
\tightlist
\item
  Define a grid of parameter values
  \(\boldsymbol\theta_1, \dots, \boldsymbol\theta_G\) covering all
  plausible values.
\item
  At each grid point compute the \textbf{log} prior plus the
  \textbf{log} likelihood (logs avoid underflow, as in Week 4).
\item
  Subtract the maximum, exponentiate, and normalize so that the values
  sum to 1 (or integrate to 1).
\end{enumerate}

The result is a discrete approximation of the posterior from which we
can compute any summary. Grid approximation is exact up to the grid
resolution and is ideal for building intuition. It fails beyond 2--3
parameters (the curse of dimensionality, Week 3), which motivates MCMC
(Weeks 11--13).

\subsubsection{5. Summarizing a posterior}

\begin{longtable}[]{@{}
  >{\raggedright\arraybackslash}p{(\linewidth - 4\tabcolsep) * \real{0.3333}}
  >{\raggedright\arraybackslash}p{(\linewidth - 4\tabcolsep) * \real{0.3333}}
  >{\raggedright\arraybackslash}p{(\linewidth - 4\tabcolsep) * \real{0.3333}}@{}}
\caption{\centering 5. Summarizing a posterior}\tabularnewline
\toprule\noalign{}
\begin{minipage}[b]{\linewidth}\raggedright
Summary
\end{minipage} & \begin{minipage}[b]{\linewidth}\raggedright
Definition
\end{minipage} & \begin{minipage}[b]{\linewidth}\raggedright
Remark
\end{minipage} \\
\midrule\noalign{}
\endfirsthead
\toprule\noalign{}
\begin{minipage}[b]{\linewidth}\raggedright
Summary
\end{minipage} & \begin{minipage}[b]{\linewidth}\raggedright
Definition
\end{minipage} & \begin{minipage}[b]{\linewidth}\raggedright
Remark
\end{minipage} \\
\midrule\noalign{}
\endhead
\bottomrule\noalign{}
\endlastfoot
Posterior mean & \(\mathbb{E}[\theta \mid \mathbf{y}]\) & Minimizes
expected squared error \\
Posterior median & 50th percentile & Robust to skew \\
MAP & \(\arg\max_\theta p(\theta\mid\mathbf{y})\) & Equals the MLE under
a flat prior \\
Posterior SD & \(\sqrt{\mathrm{Var}[\theta\mid\mathbf{y}]}\) & The
Bayesian counterpart of a SE \\
95\% equal-tailed interval & 2.5th to 97.5th percentiles & Easy to
compute \\
95\% \textbf{HDI} (highest-density interval) & The narrowest interval
containing 95\% of the posterior; every value inside is more probable
than every value outside & Preferred for skewed posteriors \\
\end{longtable}

\textbf{Credible vs.~confidence interval.} A 95\% credible interval
\emph{does} mean ``given the model, the prior and the data, the
parameter lies in this interval with probability 0.95''. A confidence
interval is a statement about the long-run behavior of the procedure
(Week 5). With weak priors and plenty of data, the two often coincide
numerically, but their interpretations differ.

\textbf{Posterior probabilities of hypotheses} come for free:
\(P(d' > 1 \mid \mathbf{y})\) is just the posterior mass above 1.

\subsubsection{6. Prior predictive checks}

The \textbf{prior predictive distribution} describes the data the model
expects \emph{before} seeing any data:

\[
p(\mathbf{y}^{\text{rep}}) = \int p(\mathbf{y}^{\text{rep}}\mid\boldsymbol\theta)\, p(\boldsymbol\theta)\, d\boldsymbol\theta.
\]

To simulate from it: draw \(\boldsymbol\theta\) from the prior, then
simulate data from the model. If the simulated data look absurd (e.g.,
nearly all hit rates at exactly 0 or 1, or recall that \emph{increases}
over a month), the prior is unreasonable. Prior predictive checks are
the Bayesian version of the ``simulate before you fit'' habit from Week
2.

\begin{quote}
\textbf{Advanced: the posterior as updated prior.} Because
\(p(\boldsymbol\theta\mid\mathbf{y}_1, \mathbf{y}_2) \propto p(\mathbf{y}_2\mid\boldsymbol\theta)\, p(\boldsymbol\theta\mid\mathbf{y}_1)\)
for independent data, today's posterior is tomorrow's prior. Bayesian
updating is therefore naturally sequential, which is also why Bayesian
inference is a popular normative model of \emph{learning} (Week 15).
\end{quote}

\paragraph{For an engineer}

\begin{itemize}
\tightlist
\item
  This is \textbf{Bayesian estimation} as in detection and estimation
  theory: the posterior is what an MAP or MMSE estimator works with. The
  posterior mean is the \textbf{MMSE estimate}; the posterior mode is
  the \textbf{MAP estimate}; with a flat prior MAP = ML.
\item
  The \textbf{conjugate Beta--binomial update} is the same idea as
  conjugate Gaussian updating, the scalar core of the \textbf{Kalman
  filter} (prior mean and variance + measurement \ensuremath{\rightarrow} posterior mean and
  variance), here for a probability instead of a Gaussian state.
\item
  \textbf{Sequential updating} (``today's posterior is tomorrow's
  prior'') is exactly recursive Bayesian filtering.
\item
  What is different in psychology: the prior is usually \textbf{weak}
  and must be \emph{justified to readers}, and the main output is not a
  decision but a description of uncertainty about a psychological
  parameter. Prior predictive checks are the habit that makes priors
  defensible.
\end{itemize}

\paragraph{For the instructor}

\begin{itemize}
\tightlist
\item
  \textbf{Do the 4-point grid on the board} (worked example, Part 1)
  with a column for prior, likelihood, product and posterior. It
  contains the whole of Bayesian inference in one table. Then change the
  prior to (0.1, 0.1, 0.4, 0.4) and recompute live: students see the
  prior at work.
\item
  \textbf{Use the Beta pseudo-count story.} ``A Beta(2, 2) prior is like
  having already seen 2 recalls and 2 failures.'' Students who find
  priors arbitrary relax once they see them as imaginary data.
\item
  \textbf{Address the ``subjectivity'' worry directly.} Every analysis
  has assumptions (the likelihood is one too); Bayesian analysis makes
  one more of them explicit, and sensitivity analysis (Week 10) checks
  how much it matters.
\item
  \textbf{Typical questions.}

  \begin{itemize}
  \tightlist
  \item
    \emph{``Which prior is correct?''} None is uniquely correct. Use
    weakly informative priors, check them with prior predictive
    simulations, and report sensitivity.
  \item
    \emph{``Is a credible interval the same as a confidence interval?''}
    Often similar numerically with weak priors and enough data, but only
    the credible interval means ``95\% probability that the parameter is
    in here''.
  \item
    \emph{``Is the brain Bayesian, then?''} That is a different question
    (Week 15). This week Bayes is a tool for the scientist.
  \end{itemize}
\end{itemize}

\begin{center}\rule{0.5\linewidth}{0.5pt}\end{center}

\Needspace{11\baselineskip}

\subsection{Worked example}

\subsubsection{Part 1: a four-point grid by hand}

A participant recalls \(y = 7\) of \(n = 10\) items. To keep the
arithmetic simple, suppose \(\theta\) can only take the values 0.2, 0.4,
0.6 or 0.8, each with prior probability 1/4.

\begin{longtable}[]{@{}
  >{\raggedright\arraybackslash}p{(\linewidth - 8\tabcolsep) * \real{0.2000}}
  >{\raggedright\arraybackslash}p{(\linewidth - 8\tabcolsep) * \real{0.2000}}
  >{\raggedright\arraybackslash}p{(\linewidth - 8\tabcolsep) * \real{0.2000}}
  >{\raggedright\arraybackslash}p{(\linewidth - 8\tabcolsep) * \real{0.2000}}
  >{\raggedright\arraybackslash}p{(\linewidth - 8\tabcolsep) * \real{0.2000}}@{}}
\caption{\centering Part 1: a four-point grid by hand}\tabularnewline
\toprule\noalign{}
\begin{minipage}[b]{\linewidth}\raggedright
\(\theta\)
\end{minipage} & \begin{minipage}[b]{\linewidth}\raggedright
Prior
\end{minipage} & \begin{minipage}[b]{\linewidth}\raggedright
Likelihood \(\theta^7(1-\theta)^3\) (constant dropped)
\end{minipage} & \begin{minipage}[b]{\linewidth}\raggedright
Prior × likelihood
\end{minipage} & \begin{minipage}[b]{\linewidth}\raggedright
Posterior
\end{minipage} \\
\midrule\noalign{}
\endfirsthead
\toprule\noalign{}
\begin{minipage}[b]{\linewidth}\raggedright
\(\theta\)
\end{minipage} & \begin{minipage}[b]{\linewidth}\raggedright
Prior
\end{minipage} & \begin{minipage}[b]{\linewidth}\raggedright
Likelihood \(\theta^7(1-\theta)^3\) (constant dropped)
\end{minipage} & \begin{minipage}[b]{\linewidth}\raggedright
Prior × likelihood
\end{minipage} & \begin{minipage}[b]{\linewidth}\raggedright
Posterior
\end{minipage} \\
\midrule\noalign{}
\endhead
\bottomrule\noalign{}
\endlastfoot
0.2 & 0.25 & 0.0000066 & 0.0000016 & \textbf{0.002} \\
0.4 & 0.25 & 0.0003539 & 0.0000885 & \textbf{0.092} \\
0.6 & 0.25 & 0.0017916 & 0.0004479 & \textbf{0.468} \\
0.8 & 0.25 & 0.0016777 & 0.0004194 & \textbf{0.438} \\
Sum & 1 & & 0.0009574 & 1.000 \\
\end{longtable}

Each posterior value = (prior × likelihood) / sum. The binomial
coefficient was dropped because it cancels in the normalization.
Interpretation: after the data, 0.6 and 0.8 are about equally plausible,
and 0.2 is practically excluded.

\subsubsection{Part 2: the same data with a continuous Beta prior}

Prior \(\mathrm{Beta}(2, 2)\) (mean 0.5, mildly favoring moderate
values). Posterior:
\(\mathrm{Beta}(2 + 7, 2 + 3) = \mathrm{Beta}(9, 5)\).

\begin{itemize}
\tightlist
\item
  Posterior mean \(= 9/14 = 0.643\): between the prior mean (0.5) and
  the observed proportion (0.7), closer to the data because
  \(n = 10 > a_0 + b_0 = 4\).
\item
  MAP (mode) \(= (9 - 1)/(9 + 5 - 2) = 0.667\).
\item
  95\% equal-tailed credible interval \(= [0.386, 0.861]\).
\end{itemize}

\textbf{Interpretation.} Given a mild prior and 10 items, the recall
probability is probably around 0.64, but anything from about 0.39 to
0.86 remains credible. Ten items tell us little.

\subsubsection{Part 3: the prior matters less with more data}

With \(y = 30\) of \(n = 50\):

\begin{longtable}[]{@{}
  >{\raggedright\arraybackslash}p{(\linewidth - 6\tabcolsep) * \real{0.2500}}
  >{\raggedright\arraybackslash}p{(\linewidth - 6\tabcolsep) * \real{0.2500}}
  >{\raggedright\arraybackslash}p{(\linewidth - 6\tabcolsep) * \real{0.2500}}
  >{\raggedright\arraybackslash}p{(\linewidth - 6\tabcolsep) * \real{0.2500}}@{}}
\caption{\centering Part 3: the prior matters less with more data}\tabularnewline
\toprule\noalign{}
\begin{minipage}[b]{\linewidth}\raggedright
Prior
\end{minipage} & \begin{minipage}[b]{\linewidth}\raggedright
Posterior
\end{minipage} & \begin{minipage}[b]{\linewidth}\raggedright
Mean
\end{minipage} & \begin{minipage}[b]{\linewidth}\raggedright
95\% credible interval
\end{minipage} \\
\midrule\noalign{}
\endfirsthead
\toprule\noalign{}
\begin{minipage}[b]{\linewidth}\raggedright
Prior
\end{minipage} & \begin{minipage}[b]{\linewidth}\raggedright
Posterior
\end{minipage} & \begin{minipage}[b]{\linewidth}\raggedright
Mean
\end{minipage} & \begin{minipage}[b]{\linewidth}\raggedright
95\% credible interval
\end{minipage} \\
\midrule\noalign{}
\endhead
\bottomrule\noalign{}
\endlastfoot
\(\mathrm{Beta}(1,1)\) (flat) & \(\mathrm{Beta}(31, 21)\) & 0.596 &
{[}0.461, 0.724{]} \\
\(\mathrm{Beta}(0.5, 0.5)\) (Jeffreys) & \(\mathrm{Beta}(30.5, 20.5)\) &
0.598 & {[}0.462, 0.727{]} \\
\(\mathrm{Beta}(20, 5)\) (informative, ``\ensuremath{\approx} 0.8'') &
\(\mathrm{Beta}(50, 25)\) & 0.667 & {[}0.557, 0.768{]} \\
\end{longtable}

The flat and Jeffreys priors give practically the same answer, also
equal to the Week 5 profile-likelihood interval {[}0.462, 0.728{]}. The
informative prior, which is worth 25 pseudo-observations, still pulls
the estimate noticeably toward 0.8: 50 observations are not enough to
overwhelm a strong prior. Whether that is good depends on whether the
prior knowledge was right.

\begin{center}\rule{0.5\linewidth}{0.5pt}\end{center}

\Needspace{11\baselineskip}

\subsection{Lab (2 hours)}

\subsubsection{Goals}

\begin{enumerate}
\def\labelenumi{\arabic{enumi}.}
\tightlist
\item
  Compute beta-binomial posteriors analytically and by grid
  approximation, and see the prior--data compromise.
\item
  Watch the posterior sharpen with data (sequential updating).
\item
  Compute a 2-D grid posterior for SDT (\(d'\), \(c\)); summarize it
  with means, HDIs and posterior probabilities.
\item
  Use prior predictive simulation to expose a bad ``vague'' prior.
\end{enumerate}

\textbf{Timing:} Part A 25 min · Part B 15 min · Part C 40 min · Part D
25 min · Part E 15 min.

\subsubsection{Part A: beta-binomial, analytically and on a grid (25
min)}

The grid and exact posteriors coincide. Look at where each posterior
sits relative to its prior and the likelihood.

\subsubsection{Part B: sequential updating (15 min)}

\textbf{What to notice:} the posterior starts as the prior, narrows as
items accumulate, and concentrates around the true value. Its width
shrinks roughly like \(1/\sqrt{n}\), as the SE did in Week 5.

\subsubsection{Part C: a two-parameter grid posterior for SDT (40 min)}

The Week 2 participant: 68 hits out of 100 signal trials and 11 false
alarms out of 100 noise trials. Priors: \(d' \sim \mathcal{N}(1, 1.5)\)
(sensitivity is usually positive, rarely above 4) and
\(c \sim \mathcal{N}(0, 1)\) (bias is usually modest).

(``1.000'' means that none of the 20 000 posterior samples fell below
the threshold, so the probability is \textgreater{} 0.9999. Report it as
such, not as certainty.)

How do the posterior mean and MAP compare with the MLE? Why are they so
close here? 2. Translate the HDI for \(d'\) and
\(P(c > 0 \mid \text{data})\) into one plain-language sentence each
about this participant. 3. Increase the prior SD on \(d'\) to 10, then
decrease it to 0.3 (centered at 1). Which summaries change, and by how
much?

\subsubsection{Part D: prior predictive checks (25 min)}

Before trusting a prior, look at the data it implies. Compare our prior
with a ``vague'' one, \(d' \sim \mathcal{N}(0, 10)\) and
\(c \sim \mathcal{N}(0, 10)\), which is sometimes used because it is
``uninformative''.

\textbf{What to notice:} the ``vague'' prior is not neutral at all. On
the probability scale it puts 96\% of simulated datasets at hit or
false-alarm rates of exactly 0 or 1, which practically never happens in
real detection experiments. It expresses a strong and absurd belief
about the data. Vagueness on the parameter scale is not vagueness on the
data scale.

Our own prior is not perfect either: about a quarter of its simulated
datasets have \(H < F\), because \(\mathcal{N}(1, 1.5)\) puts about 25\%
of its mass on negative \(d'\). For a task where participants are known
to be above chance, a prior restricted to \(d' > 0\) (e.g., a
half-normal or a Gamma) would be more defensible. Try it, and check
whether it changes the posterior for this participant. (Hint: with 200
informative trials, hardly at all.)

\subsubsection{Part E: interpret (15 min)}

Write a short ``results paragraph'' (4--6 sentences) for this
participant as it would appear in a paper. Report the priors, the
posterior means and 95\% HDIs for \(d'\) and \(c\),
\(P(d' > 1 \mid \text{data})\), and one sentence on prior sensitivity
(from Part C, question 3).

\subsubsection{Exercises (increasing difficulty)}

\textbf{Exercise 1 \ensuremath{\star}.} Show algebraically that the posterior mean of the
beta-binomial model is a weighted average of the prior mean and \(y/n\).
What are the weights?

\textbf{Exercise 2 \ensuremath{\star}.} Using the Part A grid, compute the 95\% HDI and
the equal-tailed interval for \(\mathrm{Beta}(1 + 2, 1 + 18)\) (2
successes out of 20). Why do they differ?

\textbf{Exercise 3 \ensuremath{\star}\ensuremath{\star} (forgetting by grid).} Compute the grid posterior
of \((A, d)\) for the power forgetting model with the Week 2 data, using
priors \(A \sim \mathrm{Beta}(8, 2)\) and
\(d \sim \mathrm{Gamma}(\text{shape } 2, \text{scale } 0.3)\). Compare
the posterior means and 95\% HDIs with the MLE and the Week 5 profile
and bootstrap intervals.

\textbf{Exercise 4 \ensuremath{\star}\ensuremath{\star} (equivalence of MAP and MLE).} In Part C, replace
both priors by flat priors over the grid. Show that the MAP equals the
MLE. Why is the posterior \emph{mean} still slightly different?

\textbf{Exercise 5 \ensuremath{\star}\ensuremath{\star}\ensuremath{\star} (Advanced: sequential vs.~batch).} For the SDT
model, update the grid posterior trial by trial (each trial's posterior
becomes the next trial's prior) through 200 simulated trials, and show
that the final posterior is identical (up to numerical error) to the
batch posterior from all trials at once. Plot the posterior SD of \(d'\)
against the number of trials.

\subsubsection{Lab check-off}

Show: (1) the three-panel prior/likelihood/posterior plot from Part A,
(2) the SDT summary table and joint posterior from Part C, and (3) the
two prior predictive plots from Part D with one sentence on what is
wrong with the vague prior.

\begin{center}\rule{0.5\linewidth}{0.5pt}\end{center}

\subsection{Connection to AI \& LLMs}

\textbf{How good is the model, really? A Bayesian answer for small
evaluations.} Researchers often test an LLM on a small set of items, for
example 25 hand-written problems. The beta-binomial model of this week
gives what such studies usually lack: a full posterior for the accuracy,
direct probability statements (``the probability that accuracy exceeds
60\%''), and a direct comparison of two models.

With 25 items, the posterior intervals are about 0.35 wide, and there is
a 19\% posterior probability that B is actually the better model.

How much do the conclusions change? - \ensuremath{\star}\ensuremath{\star} The comparison above treats the
two models as independent, although they answered the same items. Why is
this a problem (see the Week 5 box)? Sketch a model that respects the
pairing. - \ensuremath{\star}\ensuremath{\star} Discussion: what prior would you use for the accuracy of a
brand-new model on a four-choice test, and why should it not put much
mass below 0.25?

\begin{center}\rule{0.5\linewidth}{0.5pt}\end{center}

\subsection{Common mistakes \& misconceptions}

\begin{enumerate}
\def\labelenumi{\arabic{enumi}.}
\tightlist
\item
  \textbf{``The prior makes Bayesian inference subjective, so the
  results are arbitrary.''} Priors are explicit assumptions, just like
  the choice of likelihood. Report them, justify them, and check
  sensitivity. With adequate data, reasonable priors lead to similar
  conclusions.
\item
  \textbf{``Flat or very wide priors are uninformative.''} They can be
  strongly informative on the data scale (Part D). Use weakly
  informative priors and check them with prior predictive simulation.
\item
  \textbf{Choosing the prior after looking at the data.} That uses the
  data twice.
\item
  \textbf{Multiplying probabilities on a grid without logs,} which
  causes underflow (Week 4). Work with log-prior + log-likelihood, then
  subtract the maximum before exponentiating.
\item
  \textbf{Too coarse or too narrow a grid.} If the posterior is cut off
  at the grid edge, all summaries are wrong. Check that the posterior is
  near zero at the edges.
\item
  \textbf{Reading a credible interval as a confidence interval, or vice
  versa.} Only the credible interval supports ``95\% probability the
  parameter is in here'' (given the model and prior).
\item
  \textbf{Reporting only the MAP.} The point of the posterior is the
  whole distribution. Report a central value \emph{and} an interval, and
  probabilities for the hypotheses you care about.
\end{enumerate}

\begin{center}\rule{0.5\linewidth}{0.5pt}\end{center}

\subsection{Check your understanding}

\textbf{1.} Write Bayes' rule for a parameter \(\theta\) and name each
term. \textbf{Answer}
\(p(\theta\mid y) = p(y\mid\theta)p(\theta)/p(y)\): posterior =
likelihood × prior / evidence (marginal likelihood). For estimation,
\(p(\theta\mid y) \propto p(y\mid\theta)p(\theta)\).

\textbf{2.} Prior \(\mathrm{Beta}(3, 3)\), data: 12 recalled out of 20.
What is the posterior, and its mean? \textbf{Answer}
\(\mathrm{Beta}(3 + 12, 3 + 8) = \mathrm{Beta}(15, 11)\), with mean
\(15/26 = 0.577\).

\textbf{3.} Why does the posterior get closer to the MLE as \(n\)
increases? \textbf{Answer} The log-likelihood is a sum over observations
and grows with \(n\), while the log-prior stays fixed. With much data
the likelihood dominates the product, and the posterior concentrates
around the likelihood's maximum.

\textbf{4.} What is the difference between a 95\% HDI and a 95\%
equal-tailed interval? When do they coincide? \textbf{Answer} The HDI is
the narrowest interval containing 95\% of the posterior mass, and every
value inside has higher density than any value outside. The equal-tailed
interval cuts 2.5\% from each tail. They coincide for symmetric unimodal
posteriors and differ for skewed ones.

\textbf{5.} Why does grid approximation become impractical for models
with many parameters? \textbf{Answer} The number of grid points grows
exponentially with the number of parameters (e.g., 100 points per
parameter means \(100^k\) evaluations). That is the curse of
dimensionality. MCMC sampling (Weeks 11--13) avoids it.

\textbf{6.} A colleague uses \(d' \sim \mathcal{N}(0, 100)\) ``to be
safe''. What would you do to convince them it is a bad idea?
\textbf{Answer} A prior predictive simulation: show that this prior
implies hit and false-alarm rates that are almost always exactly 0 or 1,
an absurd belief about real detection data. Suggest a weakly informative
prior instead.

\textbf{7.} Explain the statement ``the log-linear correction for
extreme rates is a Bayesian estimate''. \textbf{Answer} The corrected
rate \((y + 0.5)/(n + 1)\) equals the posterior mean of \(\theta\) under
a \(\mathrm{Beta}(0.5, 0.5)\) prior: \((0.5 + y)/(0.5 + 0.5 + n)\). The
correction implicitly adds half a pseudo-observation of each type.

\begin{center}\rule{0.5\linewidth}{0.5pt}\end{center}

\subsection{Assignment}

\begin{itemize}
\tightlist
\item
  \textbf{A3 (model comparison and recovery)} and the \textbf{project
  proposal} are due at the start of this week's lecture.
\item
  No new graded assignment this week. \textbf{A4 (Bayesian analysis by
  grid)} is released in Week 10.
\item
  \textbf{Optional practice (ungraded):} lab Exercise 3 (forgetting by
  grid) is the core of A4. Starting it now will save you time.
\end{itemize}

\begin{center}\rule{0.5\linewidth}{0.5pt}\end{center}

\subsection{Further reading}

\begin{itemize}
\tightlist
\item
  \textbf{Kruschke (2014), chapters 5--6.} Bayes' rule, and inferring a
  binomial probability via exact mathematical analysis (the
  beta-binomial model). Very gentle and visual. \emph{(Main reading.)}
\item
  \textbf{Lee \& Wagenmakers (2014), chapters 1 and 3.} The basics of
  Bayesian analysis, and inferences with binomials.
\item
  \textbf{Farrell \& Lewandowsky (2018), chapter 6.} Bayesian parameter
  estimation: basic concepts, linked to the likelihood framework of Week
  4.
\item
  \textbf{Azimifar \& Ahmadi (1396)}, vol.~1 (in Persian): the section
  on Bayesian parameter estimation and its relation to ML estimation
  {[}verify chapter{]}.
\end{itemize}

\emph{Preview for Week 10:} the posterior tells us about parameters
\emph{within} a model. Next week we ask how the data change our belief
\emph{between} models, using the evidence
\(p(\mathbf{y}\mid\mathcal{M})\) that we have been ignoring as a
normalizing constant, and how to check whether a model's posterior
predictions look like the real data.
